% Paper 1, §10 — Multi-token output: speculation, MTP, and parallel decoding
% Anchors: kb/notes/architecture/multi-token-prediction.md
%          kb/notes/inference/speculative-decoding.md
%          kb/excerpts/{gloeckle2024-mtp,leviathan2023,cai2024-medusa,li2024-eagle}.md

\section{Multi-token output: speculation, MTP, and parallel decoding}
\label{sec:multi-token-output}
\label{sec:multitoken}

The architectural axes catalogued in
\cref{sec:tokenization-embedding,sec:positional-encoding,sec:kv-sharing,sec:hardware-attention,sec:ffn,sec:normalization-axis,sec:state-space}
all assume the canonical \glsterm{autoregressive}{autoregressive} decoding loop: one forward pass
produces one token. That assumption is load-bearing for both training
mathematics (the next-token cross-entropy of
\cref{sec:tokenization-embedding}) and serving cost (the per-step KV /
weight read of \cref{sec:kv-sharing}). The 2023--2026 literature breaks
this one-pass-one-token coupling along two distinct axes that the field
sometimes conflates. \emph{Speculative decoding}
\citep{leviathan2023,chen2023-speculative} is a serving-side algorithm
that produces multiple tokens per target-model forward pass while
preserving the target's output distribution exactly. \emph{Multi-token
prediction} (\glsterm{multi-token-prediction}{MTP}) \citep{gloeckle2024-mtp} is an architectural and
training-objective choice that adds heads predicting tokens at offsets
$t+1, \ldots, t+n$, reporting both quality lift on code-heavy
benchmarks and (when the heads are reused at inference) a built-in
self-speculative drafter. The thesis of this section is the distinction:
speculative decoding is \emph{inference-only} and never touches the
training loss; \glsterm{multi-token-prediction}{MTP} is \emph{architectural} and changes both the loss
landscape and what the deployable system can do. They compose --- a
\glsterm{precision-pinning}{DeepSeek-V3}-style \glsterm{multi-token-prediction}{MTP}-trained model \citep[\S 2.3.1]{deepseek-v3} is
its own draft model in a \glsterm{speedup-theorem}{Leviathan}-style verification loop --- but
they are not interchangeable.

\subsection{Speculative decoding: the rejection rule}
\label{sec:mt-leviathan}

Let $M_p$ be a target \glsterm{llm}{LLM} with conditional distribution $p(x_t \mid
x_{<t})$ over a \glsterm{vocabulary}{vocabulary} of size $V$, and $M_q$ a smaller draft
model with $q(x_t \mid x_{<t})$. \citet{leviathan2023} sample $\gamma$
draft tokens $x_1, \ldots, x_\gamma$ from $M_q$ autoregressively, run
$M_p$ in a single parallel forward over the $\gamma + 1$ extended
prefixes, then accept or reject each draft sequentially using
\citep[\S 2.3]{leviathan2023}\footnote{KB excerpt:
\texttt{kb/excerpts/leviathan2023.md\#sec-2-3}.}:
\begin{equation}
\text{accept } x_i \text{ with prob } \min\!\left(1, \frac{p_i(x_i)}{q_i(x_i)}\right),
\quad \text{else resample } x_i \sim \operatorname{norm}(\max(0, p_i - q_i)).
\label{eq:spec-reject}
\end{equation}
The proof in \citet[App.~A.1]{leviathan2023} shows that any token
emitted by \cref{eq:spec-reject} is i.i.d.\ from $p$ \emph{regardless
of $M_q$}. The output distribution is not approximated; it is
preserved. This is what distinguishes speculative decoding from earlier
parallel-decoding heuristics (block-parallel decoding, soft-decoding
ensembles), which trade quality for speed.

Define the per-token acceptance rate
$\alpha = \E_{x \sim q}[\min(1, p(x)/q(x))]$. Under the i.i.d.\
assumption, the expected number of tokens emitted per algorithm call
is the capped geometric \citep[\S 3.1, Eq.~1]{leviathan2023}\footnote{KB
excerpt: \texttt{kb/excerpts/leviathan2023.md\#sec-3-1}.}
\begin{equation}
\E[\#\text{tokens}] = \frac{1 - \alpha^{\gamma+1}}{1 - \alpha},
\label{eq:spec-tokens}
\end{equation}
and the expected walltime improvement, with cost coefficient $c$ equal
to the ratio of $M_q$'s per-step cost to $M_p$'s, is
\citep[\S 3.3, Thm.~3.8]{leviathan2023}
\begin{equation}
\text{walltime gain} = \frac{1 - \alpha^{\gamma+1}}{(1-\alpha)(\gamma c + 1)}.
\label{eq:spec-walltime}
\end{equation}
A \glsterm{key}{key} analytic result is \citet[Thm.~3.5, Cor.~3.6]{leviathan2023}:
$\alpha = \E[\min(p, q)] = 1 - \E[D_{LK}(p, q)]$, where $D_{LK}$ is
total-variation distance. The speedup curve is therefore a function of
how close $q$ is to $p$ in TV. Reported $\alpha$ on \glsterm{t5}{T5}-XXL with a
\glsterm{t5}{T5}-Small drafter is $\sim 0.75$, yielding $2$--$3\times$ wallclock
speedup at $\gamma = 7$ \citep[\S 3.6, Tab.~3]{leviathan2023}.
Concurrent work \citep{chen2023-speculative} derives the same rule
and is treated as joint prior in the field.

\subsection{Self-speculation without a drafter: Medusa}
\label{sec:mt-medusa}

The drafter-design problem is operational: a separate $M_q$ requires
its own training, versioning, and quantization, and a too-small drafter
has low $\alpha$ while a too-large drafter has high $c$. Medusa
\citep{cai2024-medusa} eliminates the auxiliary model by appending
$K$ small \emph{decoding heads} to the target backbone's last hidden
state $h_t$. The $k$-th head predicts the token at position $t + k +
1$ \citep[\S 2.1.1]{cai2024-medusa}\footnote{KB excerpt:
\texttt{kb/excerpts/cai2024-medusa.md\#sec-2-1-1}.}:
\begin{equation}
p_t^{(k)} = \softmax\!\left(W_2^{(k)} \cdot \bigl(\operatorname{SiLU}(W_1^{(k)} h_t) + h_t\bigr)\right),
\qquad W_1^{(k)} \in \R^{D \times D}, \; W_2^{(k)} \in \R^{D \times V},
\label{eq:medusa-head}
\end{equation}
with $W_2^{(k)}$ initialized identically to the original LM head and
$W_1^{(k)}$ initialized to zero, so each head starts as a pass-through
of the unmodified next-token distribution. To exploit multiple drafts
per step, Medusa proposes a \emph{tree} of candidates: each head emits
its top-$s_k$ tokens, and candidates are formed by Cartesian product
over heads \citep[\S 2.1.2]{cai2024-medusa}\footnote{KB excerpt:
\texttt{kb/excerpts/cai2024-medusa.md\#sec-2-1-2}.}, evaluating
$\sum_k \prod_i s_i$ continuations in one parallel target forward.
\emph{\glsterm{tree-attention}{Tree attention}} is the mechanism that makes this a single pass:
the attention mask is constructed so each candidate token attends only
to its unique ancestor chain in the tree, not to its siblings, which
shares the prefix's KV state across all branches.

Two training regimes:
Medusa-1 freezes the backbone and trains heads with the weighted
cross-entropy
$\mathcal{L}_{\text{Medusa-1}} = \sum_{k=1}^{K} -\lambda_k \log p_t^{(k)}(y_{t+k+1})$
with $\lambda_k = 0.8^k$ \citep[\S 2.2.1, Eq.~1]{cai2024-medusa};
the original LM head is unchanged, so the model's next-token
distribution is preserved exactly under \cref{eq:spec-reject}.
Medusa-2 unfreezes the backbone with the joint loss
$\mathcal{L}_{\text{LM}} + \lambda_0 \mathcal{L}_{\text{Medusa-1}}$
\citep[\S 2.2.2]{cai2024-medusa}, recovering higher acceptance at the
cost of a backbone fine-tune. Reported speedups: Medusa-1
$\sim 2.2\times$, Medusa-2 $2.3$--$2.8\times$ on Vicuna 7B/13B/33B
\citep[abstract]{cai2024-medusa}.

\citet[\S 2.3.1]{cai2024-medusa} also proposes \emph{typical
acceptance} as an alternative to \cref{eq:spec-reject}, accepting any
draft whose target probability exceeds an entropy-adaptive threshold
$\min(\epsilon, \delta \exp(-H(p)))$. This is faster but \emph{not}
distribution-preserving for temperature $\neq 0$ --- a tradeoff
sometimes elided in deployment notes.

\begin{contradiction}
Medusa with \glsterm{rejection-sampling}{rejection sampling} (\cref{eq:spec-reject} applied to
heads-as-drafter) is mathematically distribution-preserving; Medusa
with \glsterm{typical-acceptance}{typical acceptance} is not. Vendor blogs and benchmark tables
sometimes report ``Medusa speedup'' without specifying which
verification rule was used, conflating the two regimes. Any deployment
that requires reproducible sampling (eval harnesses, structured-output
RAG, \glsterm{rlhf}{RLHF} rollouts) must verify the verification rule, not just the
draft mechanism.
\end{contradiction}

\subsection{Self-speculation at the feature level: EAGLE}
\label{sec:mt-eagle}

EAGLE \citep{li2024-eagle} reframes the drafter as predicting
\emph{features} --- the second-to-top-layer \glsterm{hidden-state}{hidden state} $f$, the
input to the LM head --- rather than tokens directly. The motivating
observation \citep[abstract]{li2024-eagle}: token sequences have high
local entropy because many tokens are sampled from a flat tail, but
\glsterm{feature}{feature} sequences are smoother and more predictable, since the LM head
absorbs the entropy. The drafter is a single decoder layer (an FC
projection plus one \glsterm{transformer}{transformer} block) taking a fused input of past
features $F_{1:i}$ and the \emph{shifted} token sequence $T_{2:i+1}$,
and predicting $\hat f_{i+1}$ \citep[\S 3.1]{li2024-eagle}\footnote{KB
excerpt: \texttt{kb/excerpts/li2024-eagle.md\#sec-3-1}.}. The shift by
one is the load-bearing trick: at the time of predicting $f_{i+1}$,
the token $t_{i+1}$ that branched the actual generation has already
been sampled and is conditioned on, which resolves the multi-mode
uncertainty that pure \glsterm{feature}{feature} autoregression cannot commit to.

The drafter is trained with a combined objective \citep[\S 3.2]{li2024-eagle}:
\begin{equation}
L = \operatorname{SmoothL1}(f_{i+1}, \hat f_{i+1}) + w_{\text{cls}} \cdot \operatorname{CE}\!\left(\softmax(\text{LM\_Head}(f_{i+1})), \softmax(\text{LM\_Head}(\hat f_{i+1}))\right),
\label{eq:eagle-loss}
\end{equation}
with $w_{\text{cls}} = 0.1$. The classification term routes the
target's frozen LM head over both true and predicted features, so the
drafter optimizes \emph{token agreement} via the LM head's projection,
not raw \glsterm{feature}{feature} L2 distance. EAGLE inherits \glsterm{speedup-theorem}{Leviathan}-style rejection
sampling at the verification step \citep[\S 3.3]{li2024-eagle}, so
distribution preservation is by construction. Reported speedups on
LLaMA2-Chat 70B at MT-bench: $3.01\times$ at temperature $0$,
$2.67\times$ at temperature $1$ \citep[Fig.~1]{li2024-eagle}. EAGLE-3
\citep{eagle3-2025} (NeurIPS 2025) abandons \glsterm{feature}{feature} prediction in
favor of direct token prediction with multi-layer \glsterm{feature}{feature} fusion,
claiming up to $6.5\times$; the lineage is open at frontier scale.

\subsection{Multi-Token Prediction (MTP) as architectural choice}
\label{sec:mt-mtp}

\citet{gloeckle2024-mtp} take the dual position: rather than bolt a
drafter on at inference, modify the \emph{training} objective so the
trunk produces hidden states predictive of multiple future tokens
simultaneously. The standard next-token loss is generalized to
\citep[\S 2, Eq.~1]{gloeckle2024-mtp}\footnote{KB excerpt:
\texttt{kb/excerpts/gloeckle2024-\glsterm{multi-token-prediction}{mtp}.md\#sec-2}.}
\begin{equation}
\mathcal{L}_n = -\sum_t \log P_\theta(x_{t+n:t+1} \mid x_{1:t}) = -\sum_t \sum_{i=1}^{n} \log P_\theta(x_{t+i} \mid x_{1:t}),
\label{eq:mtp-loss}
\end{equation}
implemented with $n$ \emph{independent} output heads $f_1, \ldots,
f_n$ on a shared trunk producing $z_t = f_s(x_{1:t}) \in \R^D$ and
$P_\theta(x_{t+i} \mid x_{1:t}) = \softmax(f_i(z_t))$
\citep[\S 3]{gloeckle2024-mtp}. The training-time overhead is under
$5\%$ for $n \in \{2, 4\}$ \citep[\S 3]{gloeckle2024-mtp}.
Crucially, the next-token head $f_1$ is unchanged at inference --- so
a model can be \glsterm{multi-token-prediction}{MTP}-trained and served as if it were a vanilla
next-token model, with the auxiliary heads either discarded or reused
as drafters. Headline quality results: 13B-parameter \glsterm{multi-token-prediction}{MTP} models solve
$12\%$ more HumanEval and $17\%$ more MBPP problems than matched
next-token baselines \citep[abstract]{gloeckle2024-mtp}; the method is
\emph{increasingly} useful at larger scale.

\glsterm{precision-pinning}{DeepSeek-V3} \citep[\S 2.3.1]{deepseek-v3} productionizes a chained
variant: instead of $n$ independent heads each conditioned on $z_t$
alone, a sequence of $k$ small \glsterm{transformer}{Transformer} blocks $M_i$ refine a
chain of hidden states using the previous predicted token's embedding,
\begin{equation}
\mathbf{h}_t^{i} = M_i\!\left(\operatorname{Norm}(\mathbf{h}_t^{i-1}) ;\, \operatorname{Norm}(E_{\text{in}}[x_{t+i}])\right), \quad
\mathrm{logits}_{t+i+1}^{(i)} = \mathbf{h}_t^{i} E_{\text{out}},
\label{eq:mtp-deepseek}
\end{equation}
with the LM head $E_{\text{out}}$ \emph{shared} across all heads
(consistent with the embedding-tying choice of
\cref{sec:tokenization-embedding}). The total objective adds an
auxiliary term per chained head with weight $\lambda$ that decays
during training:
$\mathcal{L} = \mathcal{L}_{\text{NTP}} + \frac{\lambda}{D} \sum_{i=1}^{k} \mathcal{L}_{\text{MTP},i}$
\citep[\S 2.3.1]{deepseek-v3}. \glsterm{precision-pinning}{DeepSeek-V3} ships $k = 1$ (one \glsterm{multi-token-prediction}{MTP}
module beyond the main head) and reports $\sim 85\%$ acceptance for
self-speculation and $\sim 1.8\times$ wall-clock speedup
\citep[\S 2.3.1]{deepseek-v3}. The parallel-vs-chained distinction
matters because chained heads condition each future-token prediction
on the previous prediction --- the natural \glsterm{autoregressive}{autoregressive} structure
\cref{eq:spec-reject} exploits at inference.

\begin{intuition}
\glsterm{multi-token-prediction}{MTP} is one model, many heads, one objective. The shared trunk produces
a single \glsterm{hidden-state}{hidden state} $z_t = f_s(x_{1:t})$; the heads $f_1, \ldots,
f_n$ each project it to a different future-position vocab
distribution. The objective in \cref{eq:mtp-loss} is the sum of cross-
entropies of those projections against the corresponding ground-truth
tokens. Returning to the canonical multi-objective form: the gradient
of $\mathcal{L}_n$ flows back through every $f_i$ into the trunk, so
the trunk's hidden states are shaped to encode information predictive
of \emph{all} of $x_{t+1}, \ldots, x_{t+n}$, not just $x_{t+1}$.
DeepSeek's chained variant in \cref{eq:mtp-deepseek} adds a small
recurrent refinement: each $M_i$ takes the previous hidden plus the
embedding of the previous predicted token, so the $i$-th head's
distribution is conditional on the realized chain rather than only on
$z_t$. Either way, the training signal is multi-objective; either way,
inference can use just $f_1$ at zero overhead.
\end{intuition}

\subsection{The training-vs-inference axis}
\label{sec:mt-axis}

The four mechanisms above sit on a spectrum indexed by \emph{when} the
multi-token capability is paid for:

\begin{itemize}
\item \textbf{Speculative decoding} \citep{leviathan2023}: pure
inference. The target model is unchanged; the drafter $M_q$ is trained
separately or pulled from the same family. Distribution-preserving by
construction. Cost: separate drafter to deploy.

\item \textbf{Medusa-1} \citep{cai2024-medusa}: post-hoc fine-tune.
The backbone is frozen; only the heads of \cref{eq:medusa-head} are
trained. Distribution-preserving under \cref{eq:spec-reject}. Cost:
one fine-tuning stage; no drafter; no backbone change.

\item \textbf{Medusa-2 / EAGLE}: post-hoc with backbone unfrozen
(Medusa-2) or with a \glsterm{feature}{feature}-level drafter trained on the target's
features (EAGLE). Distribution-preserving under inherited rejection
sampling. Cost: backbone fine-tune (Medusa-2) or drafter \glsterm{parameters}{parameters}
$\sim 1\text{--}3\%$ of target (EAGLE).

\item \textbf{\glsterm{multi-token-prediction}{MTP}} \citep{gloeckle2024-mtp,deepseek-v3}: pretraining-
time architectural choice. The training objective is
\cref{eq:mtp-loss} (parallel) or \cref{eq:mtp-deepseek} (chained).
The trunk's hidden states are shaped by the multi-token signal during
pretraining; the heads are integrated into the architecture, not
bolted on. Distribution-preserving when reused as drafter under
\cref{eq:spec-reject}. Cost: $\sim 5\%$ training-time overhead and
$\sim 1$ \glsterm{transformer}{Transformer} block of \glsterm{parameters}{parameters} (\glsterm{precision-pinning}{DeepSeek-V3}).
\end{itemize}

The \glsterm{key}{key} invariant: every mechanism above commutes with
\cref{eq:spec-reject}, so they can be \emph{composed}. An \glsterm{multi-token-prediction}{MTP}-trained
model can use its built-in heads as the drafter, dropping the cost
coefficient $c$ in \cref{eq:spec-walltime} to nearly zero (no separate
forward pass) and pushing $\alpha$ high (the drafter's distribution
\emph{is} the target's marginal distribution at offset $i$, so
$D_{LK}(p, q)$ is small by construction). The compositionality is why
\glsterm{multi-token-prediction}{MTP} is increasingly the convergent choice at frontier 2026: it
collapses two engineering investments (drafter training; drafter
maintenance) into one (a slightly modified pretraining objective)
while leaving \cref{sec:kv-sharing,sec:hardware-attention,sec:ffn,sec:state-space}
untouched. Whether the quality lift Gloeckle reports survives at
non-DeepSeek frontier scale, and how \glsterm{multi-token-prediction}{MTP} interacts with the long
\glsterm{chain-of-thought}{chain-of-thought} generation patterns of reasoning models like
DeepSeek-R1, are open --- and the subject of the contradiction
\citet{gloeckle2024-mtp} surface vs.\ \citet{deepseek-v3} surface
about \emph{whether \glsterm{multi-token-prediction}{MTP} is mainly a speedup or also a quality lift},
addressed in \cref{sec:contradictions}.

\subsection{Forward link}
\label{sec:mt-forward}

The mechanisms in this section assume a token-stream output: the
target distribution is over a single discrete \glsterm{vocabulary}{vocabulary} $\{1, \ldots,
V\}$, and ``multi-token output'' means more positions in that same
stream. \cref{sec:multimodal} relaxes this assumption in a different
direction: the input and output streams may include image patches,
audio frames, or video tokens that are not natively part of $V$. The
three-piece pattern there --- vision encoder, projector, \glsterm{llm}{LLM} body ---
turns out to be largely orthogonal to the choices in this section: a
multimodal \glsterm{llm}{LLM} with \glsterm{multi-token-prediction}{MTP} heads is a thing one can build (and \glsterm{precision-pinning}{DeepSeek-V3}
implicitly is, in its multimodal-extended variants), and a
speculatively-decoded multimodal model is also a thing. The two axes
multiply rather than interact.
